% MAS1004 Assignment 2: long report.
%
% You fill this in with your agent, in your repository. AGENTS.md tells the
% agent how. Red text in the PDF is what still has to be replaced.
%
% - Pictures go in figures/. Copy them from results/, which is not in git.
% - Outputs go in outputs/: copy experiments.csv, LOG.md and the outputs of
%   check_submission.py and curate.py count here. \outputfile{outputs/...}
%   prints a file exactly as it is. Every number in this report has to come
%   from such an output.
%
% To make the PDF: compress this folder into a zip, and on Overleaf choose
% New Project, Upload Project. Set Menu, Compiler to XeLaTeX, and Recompile.

\documentclass[11pt]{article}
\usepackage{mas1004}

\setlength{\emergencystretch}{2em}

\reporttitle{Assignment 2: Long report}
\studentname{\fillin{your name}}
\studentid{\fillin{your student ID}}

\begin{document}
\makeheader

RomanJacketMan, last releases
\texttt{track-a-1} (``starter recipe'') and \texttt{track-b-2} (``starter model,
8 epochs'').

% The sections below are a suggestion. Change, merge or add sections as you
% like; README.md, "The long report", says what a reader should be able to see.

% ---------------------------------------------------------------------------
\section*{Track A}

\begin{center}\footnotesize
\begin{tabular}{lp{0.58\linewidth}r}
\toprule
tag & what changed from the run before & validation \\
\midrule
baseline, seed 0 & the starter recipe unchanged; all 44{,}257 training photos;
  no \texttt{selection.csv} & 0.5027 \\
baseline, seed 1 & only the seed differs from baseline seed 0 & 0.4941 \\
smooth03 & \texttt{label\_smoothing} in \texttt{loss\_fn} 0.1 $\to$ 0.3 & 0.4932 \\
crop07 & crop \texttt{scale} in \texttt{train\_transform} 0.35--1.0 $\to$
  0.7--1.0; label smoothing still 0.3 & 0.4927 \\
\bottomrule
\end{tabular}
\end{center}

\textbf{Why, in my words.}
\begin{itemize}
\item baseline: \emph{같은 설정에서도 무작위성에 의한 점수 변화가 얼마나 되는지
  확인하려고 두 번 돌렸다.} (I ran it twice to see how much the score moves
  because of randomness alone, with the settings unchanged.)
\item smooth03: \emph{라벨 스무딩이 되어있었음에도 자신있게 틀린 오답이 있었으니
  라벨 스무딩 값을 바꾼다. 개선 기준은 0.01이다.}
\item crop07: \emph{구조물이나 사물이 너무 메인이여서 생기는 문제를 방지하기 위해
  크롭을 약하게 한다 (크롭 하한 0.35 $\to$ 0.7). 기준선 대비 0.01 이상 오르지
  않으면 개선되었다고 말하기 어렵다.}
\end{itemize}

\textbf{What I think each result means.}
\begin{itemize}
\item baseline: \emph{운에 의해 변동된 점수가 그 정도 차이이니 앞으로 레시피를
  바꿔서 테스트해 봤을 때 점수 변동이 그 이하면 함부로 개선되었다고 보기
  어렵다는 뜻이다.} (That is how much the score moves by luck, so from now on,
  when I test a changed recipe, a change of less than that is not enough to call
  it an improvement.)
\item smooth03: \emph{이전에 나온 시드간 변동을 반올림해 정해둔 기준인 0.01 이상
  결과가 바뀌지 않았으니 개선되었다고 말하기 어렵다.}
\item crop07: \emph{이번에도 직전에 라벨 스무딩을 조정한 모델보다 검증점수 차이가
  -0.0005 밖에 나지 않았다. 앞서 정한 기준 0.01보다 작은 숫자이기 때문에 모델이
  개선되었다고 보기 어렵다.}
\end{itemize}

All four runs used the whole selection and the official budget: 249{,}984 photos,
batch 64, 3{,}906 steps. baseline seed 0 and seed 1 and smooth03 ran on a Colab
T4 (2{,}253.2 s, 2{,}226.4 s and 2{,}280.7 s); crop07 ran on the CPU
(10{,}887.5 s).

My last Track A release is \texttt{track-a-1}, titled ``starter recipe'': the
recipe as the assignment gave it, with all 44{,}257 training photos and no
\texttt{selection.csv}. Its validation score was 0.5027 with seed 0 and 0.4941
with seed 1. Its leaderboard score is not in yet.

The recipe in \texttt{track\_a/recipe.py} is crop07, and I did not release it:
its validation score, 0.4927, is not above the starter recipe's 0.5027, so
\texttt{track-a-1} stays my Track A entry.
\emph{개선 기준인 +0.01 점보다 점수 변화가 작았기 때문에 crop07은 개선된 모델이라고
보기 어렵기 때문에 track-a-1을 최종으로 뒀다.}

% ---------------------------------------------------------------------------
\section*{Decisions you made from photos}

\picturefile{figures/mistakes_baseline_seed0_busan_00.png}

\textbf{What I looked at.} The mistakes of \texttt{baseline} seed 0: the 48
validation photos per city the model was most confident about and still got
wrong, from \texttt{python src/curate.py mistakes
results/baseline\_seed0/model.onnx --every 2}, printed as one sheet per city in
\texttt{figures/}. The run printed, per city, how many validation photos it got
wrong: seoul 516 of 1218, busan 435 of 695, tokyo 687 of 1940, osaka 702 of
1064, taipei 493 of 1122.

\textbf{What I saw in common, in my words.}
\emph{공통점 1. 도시의 전체적인 모습보다 사물, 구조물에 더 포커스가 맞춰진
사진이다. 공통점 2. 시장, 지하철역, 공원 같이 세계 어디서나 볼 수 있는
장소들이다.}

\textbf{The rule I wrote before acting.} \emph{없었다.}

\textbf{What I changed.} There is no \texttt{track\_a/selection.csv} and no
\texttt{track\_a/dropped.csv}, and \texttt{curate.py count} prints all 44{,}257
starter training photos selected, so the scores before and after are the
baseline's 0.5027 and 0.4941. \emph{시트를 본 뒤 공통점을 본 결과 도시의 전체적인
모습보다 사물 하나에 포커스가 된 사진이 많은 것 같아 크롭 실험을 진행하게
되었다. 하지만 실험 결과가 개선 되었다고 볼 수 있는 기준인 +0.01점 보다 적은
점수 변동이었어서 track-a-1을 최종으로 뒀다.}

\textbf{The requests I gave my agent.} Three requests, as I typed them:

\begin{quote}
\noindent 오답 시트를 만들 때:\\
Please prepare two things for Colab in one queue: (1) the Step 3 mistakes
sheets for the baseline model (seed 0), and (2) the first Track B model with
the starter train\_b.py, unchanged. Tell me how long each takes on Colab, then
write colab\_runs.txt, push it, and tell me what to do in Colab. Also count the
training photos in data/starter/train and show me the number.

\noindent 크롭 설정을 확인할 때:\\
Show me exactly how track\_a/recipe.py currently crops and augments the photos
(which functions, which random crop scale range, flips, color changes, and
their current values). Do not change anything.

\noindent 크롭 변경을 요청할 때(마지막 버전):\\
Change only the scale lower bound in T.RandomResizedCrop in train\_transform()
of track\_a/recipe.py from 0.35 to 0.7, so the scale becomes (0.7, 1.0). Change
nothing else. First write my reasoning into LOG.md exactly as I give it, in
Korean, without editing or translating, under a new entry for this run. Then
show me the git diff with one sentence about the change, run -{}-quick, prepare
one Colab run with seed 0 only under a new tag, push, and tell me what to do in
Colab.
\end{quote}

\outputfile{outputs/curate_count.txt}

% ---------------------------------------------------------------------------
\section*{Track B}

\textbf{What I tried.} Two steps. The first was \texttt{src/train\_b.py} as it
came --- ResNet18 at 224 pixels from the \texttt{IMAGENET1K\_V1} weights, 4
epochs over \texttt{data/starter/train}, 176{,}896 photos through the model, on
a Colab T4 (1{,}565.9 s): validation 0.484 (seoul 0.552, busan 0.382, tokyo
0.609, osaka 0.377, taipei 0.499). \emph{Track B에서 나중에 설정을 바꿨을 때
좋아졌는지 비교할 기준이 필요했다. 그래서 처음에는 \texttt{train\_b.py}의 기본
설정인 ResNet18, 224px, 4에폭으로 starter 학습 사진만 사용해 돌렸다.}

\textbf{What I think it means.} \emph{전체 검증 점수는 0.484였는데, 도시마다
차이가 컸다. 도쿄는 0.609였지만 부산은 0.382, 오사카는 0.377이었다. 다음
실험에서는 전체 점수가 오르는지와 함께, 특히 부산과 오사카를 더 잘 맞히는지도
보고 싶다.}

\textbf{The second step, \texttt{track-b-2}.} The same script with one number
changed: the \texttt{-{}-epochs} default in \texttt{src/train\_b.py}, from 4 to
8. Nothing else differs --- the same weights, the same training photos, the same
augmentation, the same batch of 64, and the one-cycle schedule is built from
\texttt{args.epochs}, so it stretches to 8 epochs by itself. 353{,}792 photos
went through the model in 3{,}124.0 s on the same Colab T4: validation 0.5041
(seoul 0.5671, busan 0.3899, tokyo 0.639, osaka 0.3891, taipei 0.5354).
\emph{에폭을 올리는 이유: 지금 모델이 학습 사진을 더 여러 번 보면, 처음 보는
사진도 더 잘 맞히지 않을까라는 생각에 에폭을 올렸다.}

That is 0.0201 above the 4-epoch run, and every city is higher: seoul +0.0151,
busan +0.0079, tokyo +0.0299, osaka +0.0118, taipei +0.0361. The training loss
fell in all eight epochs, from 1.378 to 0.691.

\textbf{What I think it means.} \emph{점수 향상이 기준치인 +0.01의 두 배라는 점,
5개 도시의 점수가 전부 올랐다는 점, 손실이 계속 하락했다는 점에서 모델이
개선되었다고 볼 수 있을 것 같다.}

\textbf{Extra photos.} I collected none: \texttt{submission/submission.json}
names \texttt{data/starter/train} as the data, so the model saw only the starter
training photos.

\textbf{Validation next to leaderboard.} Validation 0.5041 for
\texttt{track-b-2} (0.484 for the 4-epoch run before it); the leaderboard
score is not in yet. \emph{검증과 리더보드가 서로 다른 사진으로  채점하기 때문이다.}

\subsection*{Your submission, as the grader sees it}
\outputfile{outputs/check_submission.txt}

\outputfile{outputs/submission.json}

% ---------------------------------------------------------------------------
\section*{What made the difference}

\textbf{What made the difference.} Measured on \texttt{data/starter/val}, no
change I tried in Track A beat the starter recipe: smooth03 was 0.0095 lower
than baseline seed 0 (0.4932 against 0.5027), and crop07 was 0.0005 lower than
smooth03 (0.4927 against 0.4932), while the two baseline seeds alone differ by
0.0086 (0.5027 and 0.4941). In Track B the one change I made did move the
score: 8 epochs instead of 4 gave 0.5041 against 0.484, which is 0.0201 higher
and higher in all five cities (seoul +0.0151, busan +0.0079, tokyo +0.0299,
osaka +0.0118, taipei +0.0361), and 0.0201 is more than the 0.0086 that
separates the two baseline seeds from each other.

\textbf{One more week.} The plan I wrote down before the 8-epoch run was:
\emph{다음에는 같은 모델과 학습 사진으로 에폭을 4에서
8로 늘려보고 싶다. 4에폭까지 학습 손실은 계속 줄었지만 검증 점수는 마지막에 한
번만 확인했다. 더 오래 학습했을 때 검증 점수도 오르는지, 그리고 내가 관심 있는
사물 중심 사진에서 오답이 줄어드는지 확인하고 싶다.}

% ---------------------------------------------------------------------------
\section*{All runs}
\outputfile{outputs/experiments.csv}

\section*{Lab notebook}
\outputfile{outputs/LOG.md}

\end{document}
